\section{Reduction to termination}

\begin{theorem}[Crossing words are first-passage]
  \label{thm:collatz-crossing-firstpassage}
  \lean{Obsmetrologia.Closure.Collatz.firstCoeffCrossing_word_is_firstPassage}
  \leanok
  \uses{thm:collatz-firstpassage-iff}
  The first coefficient crossing of a seed --- first contracting prefix length, all shorter
  parity words noncontracting --- yields a first-passage word.
\end{theorem}
\begin{theorem}[Descent below self reduces to termination]
  \label{thm:collatz-reduction}
  \lean{Obsmetrologia.Closure.Collatz.reachesTerminal_of_all_fall_below}
  \leanok
  \uses{thm:collatz-descent-iff}
  If every seed above the terminal cycle eventually drops below itself, every positive seed
  reaches the terminal cycle (strong induction, orbit splicing). This is the proven half of the
  strict-finite reduction.
\end{theorem}
\begin{theorem}[Conditional Collatz termination]
  \label{thm:collatz-conditional}
  \lean{Obsmetrologia.Closure.Collatz.collatz_from_parentBarrierGap_and_crossing}
  \leanok
  \uses{thm:collatz-all-firstpassage-descent, thm:collatz-crossing-firstpassage, thm:collatz-reduction}
  The gap together with the crossing hypothesis imply accelerated Collatz termination. This
  conditional depends only on \texttt{propext} and \texttt{Quot.sound} (axiom-clean); the
  unconditional statement (\cref{thm:collatz-accelerated}) rests on exactly the two frontier
  holes.
\end{theorem}
\begin{theorem}[Residue-gap conditional Collatz termination]
  \label{thm:collatz-residue-conditional}
  \lean{Obsmetrologia.Closure.Collatz.collatz_from_parentBarrierResidueGap_and_crossing}
  \leanok
  \uses{thm:collatz-residue-gap-reduces, thm:collatz-conditional}
  The sharpened residue inequality \texttt{ParentBarrierResidueGap}, together with the
  crossing hypothesis, implies accelerated Collatz termination. This is the proved route from
  the modular inequality $C<D\cdot r$ to termination, conditional only on crossing and on the
  universal residue inequality itself.
\end{theorem}
\begin{theorem}[Residue gap plus no infinite barrier stream]
  \label{thm:collatz-residue-no-infinite-stream}
  \lean{Obsmetrologia.Closure.Collatz.collatz_from_parentBarrierResidueGap_and_noPositiveInfiniteBarrierStream}
  \leanok
  \uses{thm:collatz-residue-conditional, thm:collatz-crossing-no-infinite-stream}
  The sharpened residue inequality together with
  \texttt{NoPositiveInfiniteBarrierStream} implies accelerated Collatz termination. This
  states both frontier inputs in their arithmetic forms: the modular
  parent-barrier inequality and absence of a positive infinite coefficient-barrier stream.
\end{theorem}
\begin{theorem}[Exact obstruction after the residue gap]
  \label{thm:collatz-residue-exact-obstruction}
  \lean{Obsmetrologia.Closure.Collatz.acceleratedCollatz_iff_noPositiveInfiniteBarrierStream_of_residueGap}
  \leanok
  \uses{thm:collatz-residue-no-infinite-stream, thm:collatz-descent-iff}
  Once the parent-barrier residue inequality is supplied, accelerated Collatz termination is
  equivalent to \texttt{NoPositiveInfiniteBarrierStream}. The forward direction does not use the
  residue gap: if every positive seed reaches \(\{1,2\}\), then every \(n>2\) has a terminal
  prefix below \(n\), hence a coefficient-contracting prefix by
  \texttt{contracting\_of\_descent}. Thus the residue gap removes the semigroup obstruction, but
  the infinite positive barrier-stream obstruction remains exact rather than cosmetic.
\end{theorem}
\begin{theorem}[Packaged Collatz frontier inputs]
  \label{thm:collatz-frontier-inputs}
  \lean{Obsmetrologia.Closure.Collatz.collatz_from_frontierInputs}
  \leanok
  \uses{thm:collatz-residue-no-infinite-stream}
  The current reduction can be read from the packaged input
  \texttt{CollatzFrontierInputs}: the modular parent-barrier residue inequality plus
  absence of positive infinite barrier streams imply accelerated Collatz termination.
  This package names the actual two-frontier state of the proof graph without adding
  a new mathematical assumption.
\end{theorem}
\begin{theorem}[Packaged frontier exactness]
  \label{thm:collatz-frontier-inputs-exact}
  \lean{Obsmetrologia.Closure.Collatz.frontierInputs_iff_residueGap_and_acceleratedCollatz}
  \leanok
  \uses{thm:collatz-frontier-inputs, thm:collatz-residue-exact-obstruction}
  The package is exact: \texttt{CollatzFrontierInputs} is equivalent to the residue
  gap together with accelerated Collatz termination. The reverse direction uses the
  proved fact that termination forces \texttt{NoPositiveInfiniteBarrierStream}. Thus
  after the modular residue gap, the no-stream input is not bookkeeping; it is exactly
  the remaining obstruction to termination.
\end{theorem}
\begin{theorem}[Frontier package collapses to the residue gap exactly at no-stream]
  \label{thm:collatz-frontier-inputs-one-gap-collapse}
  \lean{Obsmetrologia.Closure.Collatz.frontierInputs_iff_residueGap_of_residueGap_implies_noStream}
  \leanok
  \uses{thm:collatz-frontier-inputs-exact, thm:collatz-one-frontier-criterion}
  The honest two-frontier package \texttt{CollatzFrontierInputs} is equivalent to
  the single modular residue gap exactly when the residue gap implies
  \texttt{NoPositiveInfiniteBarrierStream}. This is the formal package-level test
  for whether the whole open problem has genuinely been isolated to
  \(C<D\cdot r\). The criterion itself is proved; the no-stream consequence from the
  residue gap is still the unproved one-lock implication.
\end{theorem}
\begin{theorem}[Frontier package collapses to the semigroup gap exactly at no-stream]
  \label{thm:collatz-frontier-inputs-parent-gap-collapse}
  \lean{Obsmetrologia.Closure.Collatz.frontierInputs_iff_parentBarrierGap_of_parentBarrierGap_implies_noStream}
  \leanok
  \uses{thm:collatz-frontier-inputs-one-gap-collapse, thm:collatz-residue-gap-reduces}
  The same package-level test at the original semigroup face:
  \texttt{CollatzFrontierInputs} is equivalent to \texttt{ParentBarrierGap} exactly
  when the semigroup gap implies \texttt{NoPositiveInfiniteBarrierStream}. The
  modular inequality \(C<D\cdot r\) is the residue sharpening of this same lock.
\end{theorem}
\begin{theorem}[Named one-lock collapse proposition]
  \label{thm:collatz-parent-one-lock-collapse}
  \lean{Obsmetrologia.Closure.Collatz.parentBarrierOneLockCollapse_iff_frontierInputs_parentBarrierGap}
  \leanok
  \uses{thm:collatz-frontier-inputs-parent-gap-collapse}
  The proposition \texttt{ParentBarrierOneLockCollapse} names exactly the missing
  implication \(\texttt{ParentBarrierGap}\to\texttt{NoPositiveInfiniteBarrierStream}\).
  Lean proves that this proposition is equivalent to collapsing the honest
  two-frontier package to \texttt{ParentBarrierGap}, and also to the corresponding
  residue-gap package collapse.
\end{theorem}
\begin{theorem}[Named one-lock supplies no-stream, crossing, and contraction]
  \label{thm:collatz-parent-one-lock-supplies-stream-frontier}
  \lean{Obsmetrologia.Closure.Collatz.noPositiveInfiniteBarrierStream_of_parentBarrierGap_and_oneLock}
  \lean{Obsmetrologia.Closure.Collatz.everyNonterminalHasCrossing_of_parentBarrierGap_and_oneLock}
  \lean{Obsmetrologia.Closure.Collatz.everyNonterminalContracts_of_parentBarrierGap_and_oneLock}
  \leanok
  \uses{thm:collatz-parent-one-lock-collapse, thm:collatz-crossing-no-infinite-stream, thm:collatz-eventual-contracts-no-infinite-stream}
  Once \texttt{ParentBarrierGap} and the named one-lock proposition are supplied,
  the no-positive-infinite-barrier-stream condition, crossing condition, and bare
  eventual-contraction condition are all proved consequences. This theorem is the
  forward-use receipt for the named collapse: it does not prove the lock, but it
  proves that the old crossing frontier is no longer an independent input after
  the named implication is assumed.
\end{theorem}
\begin{theorem}[Named one-lock supplies descent and termination]
  \label{thm:collatz-parent-one-lock-supplies-descent}
  \lean{Obsmetrologia.Closure.Collatz.allFallBelow_of_parentBarrierGap_and_oneLock}
  \lean{Obsmetrologia.Closure.Collatz.collatz_from_parentBarrierGap_and_oneLock}
  \lean{Obsmetrologia.Closure.Collatz.allFallBelow_of_residueGap_and_oneLock}
  \lean{Obsmetrologia.Closure.Collatz.collatz_from_residueGap_and_oneLock}
  \leanok
  \uses{thm:collatz-parent-one-lock-supplies-stream-frontier, thm:collatz-residue-gap-reduces, thm:collatz-conditional}
  With the named one-lock proposition in hand, the parent-barrier gap alone
  supplies universal descent below self and accelerated termination; the same
  result is exposed at the modular residue face \(C<D\cdot r\). This is the
  descent argument routed through the named one-lock hypothesis.
\end{theorem}
\begin{theorem}[Frontier package is gap plus named one-lock]
  \label{thm:collatz-frontier-inputs-gap-and-one-lock}
  \lean{Obsmetrologia.Closure.Collatz.frontierInputs_iff_parentBarrierGap_and_oneLock}
  \lean{Obsmetrologia.Closure.Collatz.frontierInputs_iff_residueGap_and_oneLock}
  \leanok
  \uses{thm:collatz-parent-one-lock-collapse, thm:collatz-parent-one-lock-supplies-stream-frontier, thm:collatz-residue-gap-reduces}
  The honest two-frontier package \texttt{CollatzFrontierInputs} is exactly the
  parent-barrier gap together with \texttt{ParentBarrierOneLockCollapse}; the
  same statement holds at the residue face. This is the precise accounting
  of the one-lock route: after choosing the modular inequality \(C<D\cdot r\) as
  the arithmetic lock, the only additional content is the named implication that
  this lock rules out positive infinite barrier streams.
\end{theorem}
\begin{theorem}[One-frontier reduction criterion]
  \label{thm:collatz-one-frontier-criterion}
  \lean{Obsmetrologia.Closure.Collatz.residueGap_alone_solves_collatz_iff_residueGap_implies_noStream}
  \leanok
  \uses{thm:collatz-frontier-inputs-exact, thm:collatz-residue-exact-obstruction}
  A reduction to the modular residue gap alone is equivalent to deriving
  \texttt{NoPositiveInfiniteBarrierStream} from that same residue gap. This is the
  formal criterion for the requested one-frontier collapse: the current Lean graph
  has not proved it, but it states exactly what such a proof must supply.
\end{theorem}
\begin{theorem}[Parent-gap one-frontier reduction criterion]
  \label{thm:collatz-parent-gap-one-frontier-criterion}
  \lean{Obsmetrologia.Closure.Collatz.parentBarrierGap_alone_solves_collatz_iff_parentBarrierGap_implies_noStream}
  \leanok
  \uses{thm:collatz-parent-one-lock-collapse, thm:collatz-conditional}
  A reduction to \texttt{ParentBarrierGap} alone is equivalent to deriving
  \texttt{NoPositiveInfiniteBarrierStream} from that same semigroup gap.
\end{theorem}
\begin{theorem}[Named one-lock termination criterion]
  \label{thm:collatz-parent-one-lock-termination}
  \lean{Obsmetrologia.Closure.Collatz.parentBarrierOneLockCollapse_iff_parentBarrierGap_alone_solves_collatz}
  \leanok
  \uses{thm:collatz-parent-one-lock-collapse, thm:collatz-parent-gap-one-frontier-criterion}
  The named collapse proposition is exactly the statement that \texttt{ParentBarrierGap}
  alone implies accelerated Collatz termination.
\end{theorem}
\begin{theorem}[One-frontier eventual-contraction criterion]
  \label{thm:collatz-one-frontier-eventual-contracts}
  \lean{Obsmetrologia.Closure.Collatz.residueGap_alone_solves_collatz_iff_residueGap_implies_everyNonterminalContracts}
  \leanok
  \uses{thm:collatz-one-frontier-criterion, thm:collatz-eventual-contracts-no-infinite-stream}
  Equivalently, a reduction to the modular residue gap alone is the same as proving that
  the residue gap forces \texttt{EveryNonterminalContracts}: every nonterminal seed has some
  finite prefix with \(3^{\mathrm{oddCount}}<2^{\mathrm{len}}\). This is the sharper target
  for a residue-refinement or inverse-thread attack on the second frontier.
\end{theorem}
\begin{theorem}[Parent-gap eventual-contraction criterion]
  \label{thm:collatz-parent-gap-eventual-contracts}
  \lean{Obsmetrologia.Closure.Collatz.parentBarrierGap_alone_solves_collatz_iff_parentBarrierGap_implies_everyNonterminalContracts}
  \leanok
  \uses{thm:collatz-parent-gap-one-frontier-criterion, thm:collatz-one-frontier-eventual-contracts}
  Equivalently, \texttt{ParentBarrierGap} alone solves accelerated Collatz exactly
  when it forces bare eventual coefficient contraction for every nonterminal seed.
\end{theorem}
